
\vspace*{.1in}
\section{Proposed Research and Methods}

% \begin{verbatim} 
%   Provide a clear research plan for a 9-month project and describe the proposed activities and methods.  
%   Include enough technical details to evaluate the impact of the proposed activity.  
%   For each activity, indicate the responsibility of the key investigator(s) and the associated budget.  
%  *** If the proposed application is building on an existing, currently funded DOE 
%      project, describe how the submitted leverages that work and is distinct from 
%     the existing funding. ***  [MENTION ASCR/NEAMS ! ]
% \end{verbatim} 
% \normalsize


As illustrated in Fig.~\ref{fig:fm1}, our goal of {\bf Phase I} is to develop a
training workflow to build AI surrogates for thermal-fluid solutions.  Initial
training sets will be based on time-independent {mean-flow} fields that are
already available or readily generated from Nek5000/RS and MFEM using direct
numerical or large-eddy simulations (DNS or LES) of turbulence in representative domains.  
The geometry for each dataset is embedded in the computational mesh, whose
nodes are associated with corresponding solution variables velocity, pressure,
and temperature.   Cases with varying geometry or parameters (e.g., Reynolds
number) are fed into an encoder, which builds the latent space. In this setting
we will also be able to take advantage of the ASCR-supported development of
differential simulation capabilities in MFEM in order to produce higher-quality
surrogates or to converge the optimization at high resolution given a surrogate
initial guess \cite{MFEM24, DMC24, Schmidt26}.    With a trained model, new
geometries can be encoded to generate surrogate solutions within seconds.  
This process has been successfully applied to LES flow data for the NASA
full-aircraft test model at different angles of attack and slat positions using
datasets of size $n \approx 50$M.  The encoding process uses PyTorch on GPUs
(several days on four to eight H200s) and inference takes about 30 seconds.  

This effort builds directly on DOE investments in high-fidelity simulation.
Nek5000/RS and MFEM developed under the Center for Efficient Exascale
Discretizations as part of DOE's Exascale Computing Project \cite{ceed}.  
Their high-order
discretizations, coupled with fast matrix-free operator evaluation and highly
scalable preconditioners gives them a competitive training edge over
traditional approaches (e.g., a 32$\times$ gain compared to 2nd-order finite
volume methods \cite{Min2024Wind,Tomboulides2024ABL}).  Nek5000/RS has been
extensively used for reactor thermal-hydraulics and is part of NRC-accepted
workflows, while MFEM has enabled important fusion applications in the DOE,
based both on magnetic and on inertial confinement approaches.  
   Both Nek5000/RS and MFEM run on NVIDIA, AMD, and Intel accelerators, so there 
   are no GPU portability issues when interfacing with the training.
The proposed work, however, is fundamentally distinct from prior ASCR
efforts.  Instead of advancing
simulation capabilities alone, we develop a generalizable foundation model that
unifies data across geometries and regimes, and we introduce a closed-loop,
AI-driven optimization framework that enables rapid design exploration and
mechanistic discovery.   

{\bf IM STILL WORKING HERE}

Figure \ref{fig:tax} illustrates the essential components of our work
plan, moving left-to-right from representation models (current simulation
data), to geometry and case generalization via a common latent space,
to optimization (ultimately) driven by agentic AI.  

\input figtax

The first component of the research plan, led by ANL, UIUC, and PSU, focuses on
the development of a {\bf geometry-aware thermal–fluid foundation model}.
Building on prior work in foundation modeling \cite{Vishwasrao2026AgenticPDE},
simulation and experimental data from wire-wrapped
bundles~\cite{Dai2025WireWrapped61Pin,kaeri37pin}, as well as related
configurations such as spacer grids, pebble beds, and twisted tubes, will be
mapped into {\rm a common latent space independent of mesh topology}. This
capability is particularly critical for reactor applications, where variations
in geometry—such as wire pitch, rod spacing, and bundle configuration—strongly
influence heat transfer and pressure losses.  The proposed model integrates
{\rm graph-based geometric representations}, {\rm variational autoencoders
(VAEs)} for compression of high-dimensional flow
fields~\cite{Vinuesa2022MLCFD,SoleraRico2024BetaVAE}, and {\rm conditional
latent diffusion models}~\cite{Vishwasrao2025DiffSPORT} to learn distributions
across operating regimes, as well as neural-operator architectures for
physics-consistent surrogate learning which we will investigate in Phase II.
Derivatives of the velocity and temperature fields are central to quantities of
interest (QOIs) such as pressure drop and heat transfer.  We will compute these
in the fine-scale training data, where derivatives can be accurately computed,
and encode them along with the primary data.  The resulting surrogate model
will enable fast, high-fidelity prediction of key reactor performance
quantities, including heat transfer coefficients, pressure drop, and
flow-induced mixing. For the cases considered, we have already established
clear agreement between the experimental data and the high-order simulation
codes, which will be used to validate the ML surrogates (e.g., 
Figs. \ref{fig:blocked} and \ref{fig:collins})

The MFEM team will contribute a {\bf targeted ICF miniapp} designed as a
high-fidelity data generator for training and validating the proposed surrogate
modeling framework. Building on existing capabilities in Lagrangian
hydrodynamics \cite{Dobrev2012}, ALE methods \cite{Tomov2018}, and scalable
H(div) radiation diffusion solvers \cite{Pazner2025}, this effort will focus on
producing physically consistent, multi-physics datasets that capture coupled
hydrodynamics–radiation behavior across parameterized geometries and operating
conditions. The miniapp will be tightly integrated into the AI workflow to
systematically generate training data, support active learning loops, and
enable evaluation of surrogate accuracy and generalization. In Phase I,
simplified ICF configurations will be used to demonstrate feasibility and
establish a robust pipeline connecting MFEM-generated data with the foundation
models, ensuring that learned representations remain consistent with the
underlying physical operators and constraints.
This effort directly complements the proposed AI architecture by providing
high-order, structure- preserving finite element data that aligns with the
project’s emphasis on physically consistent latent spaces. MFEM’s support for
complex, unstructured meshes and high-order discretizations enables systematic
variation of geometry and resolution, which is critical for training
geometry-aware models such as GNN-based encoders. Moreover, the compatibility
of MFEM operators with differentiable programming frameworks enables a natural
pathway to incorporate operator-level sensitivities into both surrogate
training and downstream optimization. This will create a strong link between the
simulation stack and the latent-space optimization strategies outlined in the
proposal, particularly for adjoint-informed refinement and hybrid optimization
workflows.

The second component, led by the University of Michigan, develops a {\bf
latent-space optimization and control framework}, targeting both steady and
transient operating regimes.  This effort will employ a hybrid strategy
combining {\rm Bayesian optimization}~\cite{Morita2022BOCFD}, {\rm
derivative-free optimization methods}, and 
{\color{red}VINUESA, remove? {\rm reinforcement learning}~\cite{Font2025DRLFlowControl}}.  
Such approaches are essential for
reactor-relevant design spaces that are high-dimensional, non-convex, and
constrained by safety and operational limits. Note that the MFEM team also
has experience with {\color{red}VINUESA: remove? reinforcement learning}
in for adaptive DOE-relevant
finite element simulations \cite{RL23}.
%
Optimization will be performed directly within the latent space of the
foundation model, enabling efficient exploration of design variables such as
wire pitch, rod arrangement, and operating conditions. The {\rm agentic AI
system} will orchestrate this process by dynamically selecting optimization
strategies, proposing new simulation queries, and balancing competing
objectives, including heat transfer enhancement, pressure loss minimization,
and safety margin preservation.
%
The framework will operate in a closed-loop workflow, in which candidate
designs generated in latent space are selectively validated using high-fidelity
simulations.  This ensures predictive accuracy while maintaining robustness
across operating conditions.
%
All team members will work on {\bf cross-domain validation and integration into the
Genesis Mission ecosystem}.



\vspace*{.1in}
\noindent {\bf Effort Breakdown:}
   {\bf ANL}-reactor TF data; model validation; device optimization \$100K;
   {\bf LLNL}-ICF data; model validation \$100K;
   {\bf U. Mich.}-model development; device optimization \$100K;
   {\bf PSU}-reactor TF data; model validation \$50K;
   {\bf UIUC}-heat-exchanger data; model validation \$50K.
{\tiny \bf CHANGE to REFLECT FULL BUDGETs}

uuuuu
